\documentclass[12pt,a4paper]{article}
\usepackage[utf8]{inputenc}
\usepackage[russian,english]{babel}
\usepackage{amsmath,amssymb,amsthm}
\usepackage{geometry}
\usepackage{hyperref}
\geometry{margin=2.5cm}

\title{\textbf{Meta-Stable Architectures}\\\Large Unified Technical Note (v1): Master Model for Volumes I--III}
\author{Symbiosis: Katia \& GPT-5 (Jippi)}
\date{2026}

\newtheorem{theorem}{Theorem}
\newtheorem{assumption}{Assumption}

\begin{document}
\maketitle

\begin{abstract}
This note provides a unified simulation-ready formulation of the trilogy. We harmonize notation, state a master deterministic-stochastic model, and collect theorem-level conditions used across Volumes I--III.
\end{abstract}

\selectlanguage{russian}
\begin{abstract}
Эта техническая записка дает единую согласованную формулировку трилогии, пригодную для симуляций. Унифицируются обозначения, задается мастер-модель (детерминированная и стохастическая), и собираются теоремы с условиями из томов I--III.
\end{abstract}
\selectlanguage{english}

\section{Unified Notation}
\begin{itemize}
    \item $a\in\mathbb R^n$: primary configuration state.
    \item $\Phi\in\mathbb R_+$: regulatory activity state.
    \item $C(a)=|\nabla V(a)|^2$: crisis intensity (gradient energy).
    \item $K\in\mathbb R_+$: accumulated structural knowledge.
    \item $I_k\in[0,1]$: subsystem isolation intensity.
\end{itemize}

\section{Master Deterministic Model}
\begin{align}
\dot a &= -g(C,K,t)\nabla V(a)-\lambda S(a,\Phi), \\
\dot\Phi &= -\kappa \Phi + u(a), \\
\dot K &= \nu C(a)-\delta_K K, \\
\dot I_k &= \xi(C_k-C_{\mathrm{crit},k})I_k(1-I_k).
\end{align}
with
\begin{equation}
C_{\mathrm{crit}}(K,t)=C_0e^{-\mu K}+\zeta\overline C_\tau(t),
\quad
\overline C_\tau(t)=\frac1\tau\int_{\max(0,t-\tau)}^t C(s)ds.
\end{equation}

\section{Master Stochastic Extension}
For the oscillatory-memory block $(x,y,H,K,\mu,S)$:
\begin{equation}
dy=(\alpha x+\mu y-\gamma y^3)dt+\sigma dW_t,
\end{equation}
with remaining coordinates evolving by deterministic drift.

\section{Core Assumptions}
\begin{assumption}
$V\in C^2$, radially unbounded, and $|\nabla V(a)|\ge c_1|a|$ outside a compact set.
\end{assumption}
\begin{assumption}
$S(a,\Phi)$ is locally Lipschitz and uniformly bounded: $|S|\le C_S$.
\end{assumption}
\begin{assumption}
$u(a)$ bounded: $|u(a)|\le U_{\max}$, and $g$ bounded below: $g\ge g_{\min}>0$.
\end{assumption}

\section{Collected Theorem Statements}
\begin{theorem}[Deterministic boundedness]
Under the core assumptions, $(a,\Phi)$ is forward complete and globally bounded.
\end{theorem}

\begin{theorem}[Crisis modulation dissipativity]
If $0\le\Theta\le1$ and $g\ge g_{\min}>0$, crisis modulation preserves dissipativity:
\begin{equation}
\dot V\le -\alpha|\nabla V|^2+\beta.
\end{equation}
\end{theorem}

\begin{theorem}[Mean-square boundedness]
If an It\^o-Lyapunov function $\mathcal V$ satisfies
\begin{equation}
\mathcal L\mathcal V\le -c_1\mathcal V+c_0+c_\sigma\sigma^2,
\end{equation}
then $\sup_{t\ge0}\mathbb E|X(t)|^2<\infty$.
\end{theorem}

\section{Simulation-Ready Recipe}
\begin{itemize}
    \item Use ODE solvers (e.g., RK45) for deterministic systems.
    \item Use Euler--Maruyama for SDE terms.
    \item Enforce constraints numerically: $S\in[0,1]$, $I_k\in[0,1]$.
    \item Track diagnostics: $V(a)$, $C(a)$, threshold crossings, and transition energy surrogates.
\end{itemize}

\section*{Relation to Trilogy}
This note is a companion document and does not replace Volumes I--III. It provides the harmonized master form for implementation and verification workflows.

\end{document}
